\documentclass[12pt]{article}
\usepackage[margin=0.9in]{geometry}
\usepackage[utf8]{inputenc}
\usepackage[T1]{fontenc}
\usepackage{hyperref}
\usepackage{enumitem}
\usepackage{longtable}
\usepackage{booktabs}
\usepackage{array}
\usepackage{parskip}
\usepackage{setspace}
\newcolumntype{P}[1]{>{\raggedright\arraybackslash}p{#1}}
\setlist[itemize]{leftmargin=*}
\setlist[enumerate]{leftmargin=*}
\onehalfspacing

\title{ED452B Part 1 Unit Plan Design Report\\Minecraft Education Coding FUNdamentals}
\author{Piter Garcia}
\date{March 19, 2026}

\begin{document}
\maketitle

\section{Purpose}
This Part 1 report documents the initial design of an inclusive Grade 4 computer science unit using Minecraft Education Coding FUNdamentals. The unit introduces algorithms, sequencing, debugging, and beginning loop thinking through a high-interest block-coding environment.

\section{Design Claim}
The core design claim is that inclusive computer science instruction must separate access barriers from computational-thinking evidence. A student may struggle with sign-in, navigation, reading directions, or executive-function demands while still understanding the coding concept.

\section{Context}
The unit is designed for Pine Brook Elementary's IGNITE Computer Science and Digital Literacy rotation. Students bring varied experience with Minecraft, keyboard and mouse control, reading stamina, language backgrounds, collaboration skills, and regulation needs.

\section{Theory}
The unit is grounded in Karten's inclusion framework, MTSS, UDL, and asset-based computer science pedagogy. Supports are designed into the lesson before students fail. Students receive multiple ways to access directions and multiple ways to demonstrate understanding.

\section{Goals}
Students will be able to:
\begin{enumerate}
\item create and test block-coded Agent algorithms;
\item plan ordered sequences before running code;
\item debug by comparing expected and actual outcomes;
\item explain one code choice, bug, or fix;
\item collaborate through structured peer roles.
\end{enumerate}

\section{Standards}
\begin{longtable}{@{}P{1.5in}P{5.1in}@{}}
\toprule
\textbf{Standard} & \textbf{Connection} \\
\midrule
NYS 4--6.CT.1 & Students develop, test, and revise block-coded programs. \\
CSTA 1A-AP-08 & Students create and follow algorithms. \\
CSTA 1A-AP-10 & Students use sequences and simple loops. \\
CSTA 1A-AP-11 & Students decompose problems into ordered steps. \\
CSTA 1A-AP-14 & Students debug errors in algorithms and programs. \\
\bottomrule
\end{longtable}

\section{Assessment Plan}
Assessment includes startup observation, planning sheets, code artifacts, debugging conferences, oral explanation, and exit tickets. The assessment design separates interface access, goal understanding, code accuracy, debugging strategy, and explanation.

\section{Lesson Sequence}
\begin{enumerate}
\item \textbf{Basic Moves:} enter Minecraft Education, open MakeCode, and move the Agent with a short sequence.
\item \textbf{Sequencing to Reach a Goal:} plan a longer ordered algorithm and test the Agent path.
\item \textbf{Debugging and Revision:} identify bugs, revise code, and explain a fix.
\item \textbf{Efficient Solutions and Reflection:} compare repeated code to loops and reflect on learning.
\end{enumerate}

\section{Target Student Supports}
\begin{longtable}{@{}P{1.6in}P{2.1in}P{2.8in}@{}}
\toprule
\textbf{Profile} & \textbf{Need} & \textbf{Support} \\
\midrule
Executive-functioning & Reads quickly or skips steps. & Stop--Read--Predict routine, checklist, teacher checkpoint. \\
CLD / language-processing & Dense text may hide understanding. & NPC read-aloud, visual vocabulary, partner reading, oral explanation. \\
Regulation / social communication & Ambiguous roles increase stress. & Driver/Navigator roles, posted agenda, stuck card. \\
\bottomrule
\end{longtable}

\section{Part 1 Reflection}
The initial design preserved the Minecraft coding sequence while adding inclusive instructional structures. The next revision needed to add implementation evidence, student-learning analysis, target-student assessment, and enactment commentary.

\end{document}
